\chapter{Differentiation and Radon--Nikodym}\label{ch:radon}

\begin{center}\itshape
A staircase can climb from $0$ to $1$ and still have derivative $0$ at almost every point.
Integrating that derivative returns $0$. The rise is real, so it must be hiding on a set of length $0$.
Which functions can be rebuilt by integrating a derivative?
\end{center}

Hardy and Littlewood met a cousin of that question while studying analytic functions. In Section~II of their 1930 paper they write that the problem ``is most easily grasped when stated in the language of cricket''. A batsman has a fixed list of scores for the season. The average he announces after an innings depends on how far back he starts the count, and he will start at the innings that makes the average largest. Their maximal inequality limits that flattery. The same bound is what lets Lebesgue recover a function from its averages over small balls.

Lebesgue asked in 1904 when $F(b)-F(a)$ equals the integral of $F'$. Existence of the derivative is too weak for that identity. In 1881 Vito Volterra built a function that is differentiable at every point and whose derivative is bounded and fails to be Riemann integrable. The devil's staircase of Chapter~\ref{ch:metric} rises from $0$ to $1$ while its derivative vanishes almost everywhere, so the integral of the derivative is $0$. Lebesgue proved that every monotone function is differentiable almost everywhere, and that the fundamental theorem holds exactly for the absolutely continuous functions.

Probability asks the same question about measures. When is there a density $f$ with $\nu(A)=\int_Af\,d\mu$? Johann Radon gave the answer on $\R^d$ in 1913. Otto Nikodym gave it on a general measure space in 1930: the density exists when $\nu$ gives no mass to the sets that $\mu$ ignores, once both measures are $\sigma$-finite.

\portrait{radon}{Johann Radon (1887--1956). Photograph, about 1920, author unknown. Public domain, via Wikimedia Commons.}

That density is a derivative of one measure with respect to the other. Conditional expectation is built from it in Theorem~\ref{thm:condexp}, and the tower property is Theorem~\ref{thm:tower}. Every probability on the line splits into atoms, a density, and a singular continuous part. The Cantor law, carried by the devil's staircase, is the third of those parts.

\section{Signed measures and the Hahn decomposition}

A finite signed measure on $(E,\mathcal A)$ is a map $\nu:\mathcal A\to\R$ with $\nu(\varnothing)=0$ and $\nu(\bigsqcup_kA_k)=\sum_k\nu(A_k)$ for every sequence of disjoint measurable sets. The series converges absolutely, because a real series that converges in every order is absolutely convergent (Volume~1). The difference of two finite measures is a finite signed measure. Jordan's decomposition, below, says that every finite signed measure arises in that way. A measurable set $P$ is \emph{positive} for $\nu$ when $\nu(A)\ge0$ for every measurable $A\subseteq P$, and \emph{negative} when $\nu(A)\le0$ for every such $A$.

\begin{lemma}\label{lem:positivepart}
Whenever $\nu(B)>0$, the set $B$ contains a positive set $P$ with $\nu(P)\ge\nu(B)$.
\end{lemma}
\begin{proof}
Pull the heaviest negative pieces out of $B$, one at a time. If $B$ is already positive, take $P=B$. If not, let $n_1$ be the smallest positive integer for which some measurable $C_1\subseteq B$ satisfies $\nu(C_1)<-\frac1{n_1}$. If $B\setminus C_1$ is still not positive, let $n_2$ be the smallest positive integer for which some $C_2\subseteq B\setminus C_1$ satisfies $\nu(C_2)<-\frac1{n_2}$, and continue. If the search stops, the set that remains is positive and we are done. Otherwise set $P=B\setminus\bigcup_kC_k$. The sets $C_k$ are disjoint, so
\[
\nu(B)=\nu(P)+\sum_k\nu(C_k).
\]
The series converges, and $\nu(C_k)<-\frac1{n_k}$, so $\sum_k\frac1{n_k}<\infty$ and $n_k\to\infty$. Each term $\nu(C_k)$ is negative, hence $\nu(P)\ge\nu(B)$. Suppose some measurable $A\subseteq P$ had $\nu(A)<0$. For large $k$ one has $n_k\ge2$ and $\nu(A)<-\frac1{n_k-1}$. That $A$ sits inside $B\setminus(C_1\cup\cdots\cup C_{k-1})$, which contradicts the choice of $n_k$ as the smallest such integer. So $P$ is positive.
\end{proof}

\begin{theorem}[Hahn decomposition]\label{thm:hahn}
If $\nu$ is a finite signed measure, the space splits as $E=P\sqcup N$ with $P$ positive and $N$ negative. If $P'\sqcup N'$ is a second such partition, then $\nu$ vanishes on every subset of the symmetric difference $P\,\triangle\,P'$.
\end{theorem}
\begin{proof}
Set $m=\sup\{\nu(Q):Q\text{ is positive}\}$. First, $m$ is finite. A finite or countable union of positive sets is positive whenever the union has finite measure on every subset: split the subset into disjoint pieces, one in each set of the union, and add the nonnegative values. If there were positive sets $R_j$ with $\nu(R_j)>j$, the unions $P_n=\bigcup_{j\le n}R_j$ would be positive, and $\nu(P_n)=\nu(R_n)+\nu(P_n\setminus R_n)\ge\nu(R_n)>n$, because $P_n\setminus R_n$ is a measurable subset of the positive set $P_n$. The union $P=\bigcup_nP_n$ would then satisfy $\nu(P)\ge\nu(P_n)>n$ for every $n$, since $P\setminus P_n$ would contribute a nonnegative amount. A real-valued set function cannot do that. So $m<\infty$. Choose positive sets $Q_j$ with $\nu(Q_j)\to m$, and set $P=\bigcup_jQ_j$. The same splitting shows that $P$ is positive, and $\nu(P)\ge\nu(Q_j)$ for every $j$, so $\nu(P)=m$. Let $N=E\setminus P$. If some measurable $B\subseteq N$ had $\nu(B)>0$, Lemma~\ref{lem:positivepart} would hand us a positive $P''\subseteq B$ with $\nu(P'')>0$, and $P\cup P''$ would be positive with measure greater than $m$. So $N$ is negative. For the uniqueness clause, $P\setminus P'=P\cap N'$ is both positive and negative, and so is $P'\setminus P=P'\cap N$. Every measurable subset of either piece has measure $0$.
\end{proof}

\begin{corollary}[Jordan decomposition]\label{cor:jordan}
Every finite signed measure splits as $\nu=\nu^+-\nu^-$, with $\nu^\pm$ finite positive measures, $\nu^+(A)=\nu(A\cap P)$ and $\nu^-(A)=-\nu(A\cap N)$, living on the disjoint sets $P$ and $N$. The measure
$|\nu|=\nu^++\nu^-$ is the \emph{total variation} of $\nu$.
\end{corollary}
\begin{proof}
Take a Hahn partition $E=P\sqcup N$. For every measurable $A$,
\[
\nu(A)=\nu(A\cap P)+\nu(A\cap N)=\nu^+(A)-\nu^-(A).
\]
If $A\subseteq P$, then $\nu(A)\ge0$, so $\nu^+$ is nonnegative. If $A\subseteq N$, then $\nu(A)\le0$, so $\nu^-$ is nonnegative. Countable additivity of $\nu^\pm$ is countable additivity of $\nu$ on the pieces $A\cap P$ and $A\cap N$. Both are finite because $\nu$ is real-valued. They charge disjoint sets, and $|\nu|(A)=\nu(A\cap P)-\nu(A\cap N)$.
\end{proof}

The picture below is the signed density $f(x)=x$ on $[-1,1]$, and $f=0$ outside that interval. The coral triangle lies under the axis on $[-1,0]$, and its area is $\int_{-1}^0 x\,dx=-\frac12$. The charcoal triangle lies over the axis on $[0,1]$, and its area is $\int_0^1 x\,dx=\frac12$. One Hahn partition takes the positive set $P=(0,1]$ and the negative set $N=(-\infty,0]$. The origin, and everything outside $[-1,1]$, carries mass $0$, so those points may join either side. The total variation charges $\frac12$ to each triangle.

\begin{figure}[ht]
\centering
\fig{ch11-rw-hahn}
\caption{The signed density $f(x)=x$ on $[-1,1]$, with $f=0$ off that interval. The charcoal triangle is the positive mass $\frac12$ on $(0,1]$. The coral triangle is the negative mass $-\frac12$ on $[-1,0)$. Every measurable piece of the charcoal set has nonnegative $\nu$-measure, and every measurable piece of the coral set has nonpositive $\nu$-measure. The origin has mass $0$, so it may sit in either set.}
\end{figure}

\section{The Radon--Nikodym theorem}

Two measures $\mu$ and $\nu$ on $(E,\mathcal A)$ are \emph{mutually singular}, written $\nu\perp\mu$, when they live on disjoint sets: $E=S\sqcup T$ with $\mu(S)=0$ and $\nu(T)=0$. The measure $\nu$ is \emph{absolutely continuous} with respect to $\mu$, written $\nu\ll\mu$, when $\mu(A)=0$ forces $\nu(A)=0$. If $\nu(A)=\int_Af\,d\mu$ for some measurable $f\ge0$, then $\nu\ll\mu$, because the integral of $f$ over a $\mu$-null set is $0$. The theorem is the converse, and it needs $\sigma$-finiteness.

\begin{think}
Take $\mu$ to be standard Gaussian, and take $\nu$ to be Gaussian of mean $1$ and variance $1$. Does some function $f$ satisfy $\nu(A)=\int_Af\,d\mu$ for every Borel set $A$? What about $\nu$ the law
of a Poisson variable?
\end{think}

For the two Gaussians the ratio of the densities works. With $\varphi(x)=(2\pi)^{-1/2}e^{-x^2/2}$,
\[
f(x)=\frac{\varphi(x-1)}{\varphi(x)}=e^{x-\frac12}.
\]
A Poisson law lives on $\N$. That set has Gaussian measure $0$ and Poisson measure $1$, so no such $f$ exists. A Poisson law and a Gaussian law are mutually singular.

\begin{theorem}[Radon--Nikodym]\label{thm:rn}
Let $\mu$ and $\nu$ be $\sigma$-finite, and assume $\nu\ll\mu$. Some measurable $f\ge0$ satisfies $\nu(A)=\int_Af\,d\mu$ for every $A\in\mathcal A$, and $f$ is unique up to a $\mu$-null set. It is written $\frac{d\nu}{d\mu}$, the
\emph{Radon--Nikodym derivative}.
\end{theorem}
\begin{proof}
\emph{Finite measures.} Let $\mathcal G$ be the class of measurable $g\ge0$ such that $\int_Ag\,d\mu\le\nu(A)$ for every measurable $A$, and set $m=\sup_{g\in\mathcal G}\int g\,d\mu$. Then $m\le\nu(E)$. If $g,h\in\mathcal G$, the maximum $\max(g,h)$ is in $\mathcal G$: split $A$ into $A\cap\{g\ge h\}$ and $A\cap\{g<h\}$, and add the two inequalities. So there is a sequence $g_n\in\mathcal G$ with $\int g_n\to m$. Replacing $g_n$ by $\max(g_1,\dots,g_n)$ makes the sequence increasing and keeps it in $\mathcal G$. Monotone convergence (Theorem~\ref{thm:mct6}) gives $f=\lim g_n\in\mathcal G$ and $\int f\,d\mu=m$. The set function $\nu_0(A)=\nu(A)-\int_Af\,d\mu$ is a finite measure, and $\nu_0\ll\mu$. The claim is $\nu_0=0$. If not, $\nu_0(E)>0$, so $\mu(E)>0$, and some $\varepsilon>0$ satisfies $\nu_0(E)>\varepsilon\mu(E)$. Let $P\sqcup N$ be a Hahn decomposition of the signed measure $\nu_0-\varepsilon\mu$. The negative set contributes a nonpositive amount, so
\[
(\nu_0-\varepsilon\mu)(P)\ge(\nu_0-\varepsilon\mu)(E)>0.
\]
In particular $\mu(P)>0$: otherwise $\nu_0(P)=0$ as well, and the display would be $0$. For every measurable $A$, positivity of $P$ gives $\varepsilon\mu(A\cap P)\le\nu_0(A\cap P)$, and $f\in\mathcal G$ gives $\int_{A\setminus P}f\,d\mu\le\nu(A\setminus P)$. Therefore
\begin{align*}
\int_A(f+\varepsilon\mathbf 1_P)\,d\mu
&=\int_Af\,d\mu+\varepsilon\mu(A\cap P)\\
&\le\int_Af\,d\mu+\nu_0(A\cap P)\\
&=\int_{A\setminus P}f\,d\mu+\nu(A\cap P)\le\nu(A).
\end{align*}
So $f+\varepsilon\mathbf 1_P\in\mathcal G$. Its integral is $m+\varepsilon\mu(P)>m$, which contradicts the choice of $m$.

\emph{$\sigma$-finite measures.} Write $E$ as a countable union of sets of finite $\mu$-measure, and as a countable union of sets of finite $\nu$-measure. The intersections of those pieces, made disjoint, are sets $E_n$ on which both measures are finite. Apply the finite case on each $E_n$ and add the densities.

\emph{Uniqueness.} Suppose $\int_Af\,d\mu=\int_Ag\,d\mu$ for every measurable $A$. On $A=\{f>g\}\cap E_n\cap\{f\le k\}$ the integrand $f-g$ is positive and bounded, and the integral over $A$ is $0$. Proposition~\ref{prop:markov}(b) forces $\mu(A)=0$. Let $k\to\infty$ and then $n\to\infty$. The set $\{g>f\}$ is null by the same argument with $f$ and $g$ exchanged.
\end{proof}

Exercise~\ref{ex:v7c11b} shows why $\sigma$-finiteness stays in the hypothesis: counting measure on $\R$ dominates Lebesgue measure and still has no density. The finite case of the proof is the one conditional expectation will call on, since a probability is a finite measure.

If $\nu=f\mu$ and $X\ge0$ is measurable, the identity $\int X\,d\nu=\int Xf\,d\mu$ holds. It is the definition of $f$ when $X=\mathbf 1_A$. It passes to nonnegative simple functions by addition (Corollary~\ref{cor:linearity6}), and to every measurable $X\ge0$ by monotone convergence. For an integrable signed $X$, split $X=X^+-X^-$. In probability this is the change of measure: if $\mathbb Q\ll\mathbb P$ and $L=\frac{d\mathbb Q}{d\mathbb P}$, then $\mathbb E_{\mathbb Q}[X]=\mathbb E_{\mathbb P}[XL]$ whenever $X$ is integrable under $\mathbb Q$, or nonnegative. Reweighting samples by $L$ is importance sampling. The same ratio is the comparison in the Neyman--Pearson lemma, and it is the density Girsanov's theorem uses in Volume~12. Those two applications sit downstream of the identity just proved. Nothing below calls on them.

\begin{theorem}[Lebesgue decomposition]\label{thm:lebdecomp}
Let $\mu,\nu$ be $\sigma$-finite. The decomposition $\nu=\nu_a+\nu_s$ is unique, with $\nu_a\ll\mu$ and $\nu_s\perp\mu$. Hence $\nu=f\mu+\nu_s$ with $f=\frac{d\nu_a}{d\mu}$.
\end{theorem}
\begin{proof}
First let $\nu$ be finite. Set $s=\sup\{\nu(Z):\mu(Z)=0\}$, and choose $\mu$-null sets $Z_j$ with $\nu(Z_j)\to s$. The union $Z=\bigcup_jZ_j$ is $\mu$-null, and $\nu(Z)=s$ because $s\ge\nu(Z)\ge\nu(Z_j)$ for every $j$. Define $\nu_s(A)=\nu(A\cap Z)$ and $\nu_a(A)=\nu(A\setminus Z)$. Then $\nu_s$ is carried by the $\mu$-null set $Z$, so $\nu_s\perp\mu$. If $\mu(A)=0$, the union $Z\cup A$ is $\mu$-null, hence $s\ge\nu(Z\cup A)=s+\nu_a(A)$, and $\nu_a(A)=0$. Thus $\nu_a\ll\mu$. Theorem~\ref{thm:rn} supplies $f$. If $\nu_a+\nu_s=\nu_a'+\nu_s'$ is another such splitting, the finite signed measure $\nu_a-\nu_a'=\nu_s'-\nu_s$ vanishes on every $\mu$-null set, because both absolutely continuous parts do. It is also carried by a $\mu$-null set, the union of the null sets that carry $\nu_s$ and $\nu_s'$. A signed measure that vanishes on null sets and lives on a null set is the zero measure: its value on $A$ equals its value on $A$ cut by that null set, which is $0$.

If $\nu$ is only $\sigma$-finite, cut $E$ into disjoint measurable pieces on which both $\mu$ and $\nu$ are finite, run the finite case on each piece, and add. A second decomposition restricts to a decomposition on each piece, so the finite uniqueness gives the global one.
\end{proof}

\begin{corollary}[the three parts of a law]\label{cor:threeparts}
Every probability $\mu$ on $\R$ is uniquely $\mu=\mu_d+\mu_{ac}+\mu_{sc}$, where $\mu_d=\sum_ip_i\delta_{a_i}$ is carried by countably many atoms, $\mu_{ac}$ has a density with respect to $\lambda$, and $\mu_{sc}$ has no atoms and
lives on a $\lambda$-null set.
\end{corollary}
\begin{proof}
The set of atoms $\{a:\mu(\{a\})>0\}$ is countable: at most $n$ of them can have mass greater than $\frac1n$. Let $\mu_d$ be the restriction of $\mu$ to that set, and let $\mu_c=\mu-\mu_d$. Then $\mu_c$ has no atoms. Theorem~\ref{thm:lebdecomp}, applied to $\mu_c$ and Lebesgue measure, splits $\mu_c$ into an absolutely continuous part and a singular part. A restriction of an atomless measure is atomless, so the singular part has no atoms either. The atoms of $\mu$ are uniquely determined as the points of positive mass, and the Lebesgue decomposition of what remains is unique.
\end{proof}

\begin{figure}[ht]
\centering
\fig{ch10-threeparts}
\caption{The distribution function of one law with all three parts. A jump of height $\frac14$ at $-1$ is an atom. The smooth rise of height $\frac12$ on $[0,1]$ is the integral of a density. The continuous climb of height $\frac14$ on $[2,3]$ is a devil's staircase, constant on a set of length $1$ inside that interval and rising on a Cantor set of length $0$.}
\label{fig:threeparts}
\end{figure}

Exercise~\ref{ex:salem} goes further: a singular continuous law can have a distribution function that is continuous and strictly increasing, with derivative $0$ almost everywhere. The climb is packed into a dense set of length $0$.

\section{The maximal function and Lebesgue's differentiation theorem}

On $\R^d$, a continuous density can be read by zooming in. If $\nu=f\lambda_d$ and $f$ is continuous at $x$, then the average of $f$ over $B(x,r)$ tends to $f(x)$ as $r\to0$. Lebesgue's theorem says the same limit holds at almost every $x$, for every locally integrable $f$. The difficulty is that every radius has to be controlled at once. That is the job of the maximal function. A measurable $f$ is \emph{locally integrable}, written $f\in L^1_{\mathrm{loc}}$, when $\int_K|f|<\infty$ for every compact $K$. Balls are open Euclidean balls. The ball $B(x,r)$ has volume $\omega_dr^d$, where $\omega_d=\lambda_d(B(0,1))$ is the volume of the unit ball from Chapter~\ref{ch:products}.

\begin{lemma}[Vitali's covering lemma]\label{lem:vitalicover}
Given finitely many open balls $B_1,\dots,B_n$, some disjoint subcollection $B_{i_1},\dots,B_{i_k}$ satisfies $\lambda_d(B_1\cup\cdots\cup B_n)\le3^d\sum_j\lambda_d(B_{i_j})$.
\end{lemma}
\begin{proof}
Select greedily. Take a ball of largest radius, throw away every ball that meets it, take a largest ball among those that remain, and repeat. A discarded ball $B'$ meets some chosen ball $B$ of radius at least as large. Let $B$ have centre $c$ and radius $r$, and let $B'$ have centre $c'$ and radius $r'\le r$. Some point $z$ lies in both balls, so $|c'-c|<r'+r\le2r$. Any $y\in B'$ then satisfies $|y-c|<r'+2r\le3r$. So $B'$ sits inside the open ball of radius $3r$ about $c$. That concentric ball has volume $3^d$ times the volume of $B$. The selected balls are disjoint, and the tripled copies cover every ball we started with, hence cover the union.
\end{proof}

On the line the factor $3$ is visible with four intervals. The coral bars in the next picture are $(0,6)$, $(5,9)$, $(10,14)$ and $(13,15)$. The longest is $(0,6)$. It meets $(5,9)$, so that interval is discarded. Among the two that remain, $(10,14)$ is longer, and it meets $(13,15)$. The charcoal bars that survive are disjoint. Stretching each of them by a factor three produces $(-6,12)$ and $(6,18)$, and those two stretches swallow $(5,9)$ and $(13,15)$.

\begin{figure}[ht]
\centering
\fig{ch11-rw-vitali}
\caption{Vitali's selection on the line. The coral bars are the open intervals $(0,6)$, $(5,9)$, $(10,14)$ and $(13,15)$. Keep a longest interval, discard every interval that meets it, and repeat. The charcoal bars that remain, $(0,6)$ and $(10,14)$, are disjoint. The lower bars are those two intervals with radius multiplied by three. Each discarded interval meets a chosen interval at least as long, so it sits inside the corresponding stretch. Length is multiplied by $3$, which is the constant $3^d$ of Lemma~\ref{lem:vitalicover} when $d=1$.}
\end{figure}

For $f\in L^1_{\mathrm{loc}}(\R^d)$, the \emph{Hardy--Littlewood maximal function} is
\[Mf(x)=\sup_{r>0}\frac1{\lambda_d(B(x,r))}\int_{B(x,r)}|f|\,d\lambda_d.\]
It records, at each $x$, the most flattering average of $|f|$ over balls centred at $x$. The set $\{Mf>t\}$ is open. If the average of $|f|$ over $B(x,r)$ exceeds $t$, continuity of the volume in the radius gives a slightly larger radius $\rho$ for which the integral over $B(x,r)$ still exceeds $t$ times the volume of $B(x,\rho)$. Every centre $x'$ with $|x'-x|<\rho-r$ then has $B(x,r)\subseteq B(x',\rho)$, so the average over $B(x',\rho)$ exceeds $t$ as well. Thus $Mf$ is Borel measurable.

\begin{theorem}[maximal inequality]\label{thm:maximal}
If $f\in L^1(\R^d)$ and $t>0$, then $\lambda_d(\{Mf>t\})\le\dfrac{3^d}t\|f\|_1$.
\end{theorem}
\begin{proof}
Take a compact set $K\subseteq\{Mf>t\}$. Each $x\in K$ is the centre of some ball $B_x$ with $\int_{B_x}|f|>t\lambda_d(B_x)$. Finitely many of those balls cover $K$. Lemma~\ref{lem:vitalicover} keeps a disjoint subcollection, and
\[
\lambda_d(K)\le3^d\sum_j\lambda_d(B_{i_j})<\frac{3^d}t\sum_j\int_{B_{i_j}}|f|\le\frac{3^d}t\|f\|_1.
\]
Take the supremum over compact $K$. Inner regularity, Proposition~\ref{prop:regularity}, turns that supremum into $\lambda_d(\{Mf>t\})$.
\end{proof}

\begin{figure}[ht]
\centering
\fig{ch10-maximal}
\caption{The indicator of $[-1,1]$, drawn in charcoal, and its maximal function, drawn in coral. The coral graph equals $1$ on the open interval $(-1,1)$, equals $\frac12$ at each endpoint, and equals $\frac1{|x|+1}$ for $|x|\ge1$. The tail $\frac1{|x|}$ keeps $Mf$ out of $L^1$, while $\lambda(\{Mf>t\})$ is still of order $\frac1t$. That is the weak bound in Theorem~\ref{thm:maximal}.}
\label{fig:maximal}
\end{figure}

\begin{theorem}[Lebesgue's differentiation theorem]\label{thm:lebdiff}
Let $f\in L^1_{\mathrm{loc}}(\R^d)$. For almost every $x$,
\[\lim_{r\to0}\frac1{\lambda_d(B(x,r))}\int_{B(x,r)}|f(y)-f(x)|\,dy=0.\]
These points $x$ are the \emph{Lebesgue points} of $f$. At each of them the averages of $f$ over $B(x,r)$ tend to $f(x)$.
\end{theorem}
\begin{proof}
The claim is local, so multiply $f$ by the indicator of a large ball and assume $f\in L^1(\R^d)$. Write
\[
\Omega f(x)=\limsup_{r\to0}\frac1{\lambda_d(B(x,r))}\int_{B(x,r)}|f(y)-f(x)|\,dy.
\]
If $g$ is continuous then $\Omega g=0$ at every point, because the integrand $|g(y)-g(x)|$ is small throughout a small ball. For $g\in C_c(\R^d)$ the triangle inequality
\[
|f(y)-f(x)|\le|(f-g)(y)|+|g(y)-g(x)|+|(g-f)(x)|
\]
gives $\Omega f\le M(f-g)+|f-g|$. Hence, for $t>0$,
\[
\lambda_d(\{\Omega f>2t\})\le\lambda_d(\{M(f-g)>t\})+\lambda_d(\{|f-g|>t\})\le\frac{3^d+1}t\|f-g\|_1.
\]
The first term is Theorem~\ref{thm:maximal}. The second is Markov's inequality, Proposition~\ref{prop:markov}(a). Theorem~\ref{thm:ccdense} says that continuous functions of compact support are dense in $L^1$, so the right-hand side can be made arbitrarily small. The outer measure of $\{\Omega f>2t\}$ is therefore $0$ for every $t>0$. The set where $\Omega f>0$ is a countable union of such null sets, and $\Omega f=0$ almost everywhere.
\end{proof}

In dimension one, the interval $[x,x+h]$ sits inside $B(x,2|h|)$ and has a quarter of that ball's length. At a Lebesgue point the average of $|f-f(x)|$ over the ball tends to $0$, so the one-sided average $\frac1h\int_x^{x+h}f$ tends to $f(x)$ as $h\to0$ from either side. If $F(x)=\int_a^xf$, then $F$ is differentiable at every Lebesgue point of $f$, with $F'(x)=f(x)$. That is one half of the fundamental theorem, now for every integrable integrand, at almost every point.

\section{Differentiating measures and monotone functions}

A Borel measure on $\R^d$ that is finite on bounded sets is called \emph{locally finite}. Differentiating such a measure uses outer regularity. The next lemma proves it for every finite Borel measure on a metric space. Chapter~\ref{ch:metricmeasures} starts from this fact. The argument here does not borrow it.

\begin{lemma}\label{lem:outerreg}
Every finite Borel measure $\nu$ on a metric space $X$ is outer regular, meaning $\nu(A)=\inf\{\nu(U):U\supseteq A\text{ open}\}$ for each Borel set $A$.
\end{lemma}
\begin{proof}
Let $\mathcal C$ be the class of Borel sets $A$ such that for every $\varepsilon>0$ there exist a closed set $C$ and an open set $U$ with $C\subseteq A\subseteq U$ and $\nu(U\setminus C)<\varepsilon$. A closed set $C$ belongs to $\mathcal C$. The open sets $U_n=\{x:d(x,C)<\frac1n\}$ decrease to $C$, and $\nu(U_1)<\infty$, so continuity from above (Proposition~\ref{prop:measureprops}(d)) gives $\nu(U_n\setminus C)\to0$. The class $\mathcal C$ is stable under complements. If $C\subseteq A\subseteq U$, then $U^c\subseteq A^c\subseteq C^c$, the set $U^c$ is closed, the set $C^c$ is open, and $C^c\setminus U^c=U\setminus C$. For a countable union $A=\bigcup_kA_k$, choose $C_k\subseteq A_k\subseteq U_k$ with $\nu(U_k\setminus C_k)<\varepsilon2^{-k-1}$. Let $U=\bigcup_kU_k$, which is open. The remainders $R_K=\bigcup_kC_k\setminus\bigcup_{k\le K}C_k$ decrease to the empty set, and $\nu$ is finite, so some finite union $C=C_1\cup\cdots\cup C_K$ is closed, sits inside $A$, and satisfies $\nu(\bigcup_kC_k\setminus C)<\varepsilon/2$. Then
\[
U\setminus C\subseteq\bigcup_k(U_k\setminus C_k)\cup\Bigl(\bigcup_kC_k\setminus C\Bigr),
\]
and subadditivity gives $\nu(U\setminus C)<\varepsilon$. So $\mathcal C$ is a $\sigma$-algebra. It contains every closed set, hence every open set, hence every Borel set. In particular, for each Borel $A$ and each $\varepsilon>0$ there is an open $U\supseteq A$ with $\nu(U\setminus A)\le\nu(U\setminus C)<\varepsilon$, which is outer regularity. The same pair $(C,U)$ is the two-sided approximation the next chapter quotes.
\end{proof}

\begin{theorem}[differentiation of measures]\label{thm:diffmeasures}
Let $\nu$ be a locally finite Borel measure on $\R^d$, and write its Lebesgue decomposition as $\nu=f\lambda_d+\nu_s$. Then for almost every $x$,
\[\lim_{r\to0}\frac{\nu(B(x,r))}{\lambda_d(B(x,r))}=f(x).\]
\end{theorem}
\begin{proof}
Local finiteness forces $f$ to be locally integrable, because $f\lambda_d$ is at most $\nu$ on each compact set. Theorem~\ref{thm:lebdiff} says that the averages of $f$ tend to $f(x)$ at almost every $x$, and those averages are exactly the ratios for $f\lambda_d$. It remains to show that $\nu_s(B(x,r))/\lambda_d(B(x,r))\to0$ for almost every $x$. The claim is local, and $\nu_s$ is finite on each ball, so we may replace $\nu_s$ by its restriction to a large ball and assume $\nu_s$ is finite. It is carried by a Borel set $Z$ with $\lambda_d(Z)=0$ and $\nu_s(\R^d\setminus Z)=0$. Fix $k\ge1$ and set
\[
F_k=\Bigl\{x\notin Z:\limsup_{r\to0}\frac{\nu_s(B(x,r))}{\lambda_d(B(x,r))}>\frac1k\Bigr\}.
\]
Given $\varepsilon>0$, Lemma~\ref{lem:outerreg} supplies an open set $U\supseteq\R^d\setminus Z$ with $\nu_s(U)<\varepsilon$. Each $x\in F_k$ lies in $U$. The limsup condition produces arbitrarily small radii, so some ball $B_x\subseteq U$ centred at $x$ satisfies $\nu_s(B_x)>\frac1k\lambda_d(B_x)$. Let $K$ be a compact subset of $\bigcup_xB_x$. Finitely many of the balls $B_x$ cover $K$, and Lemma~\ref{lem:vitalicover} keeps a disjoint subcollection with
\[
\lambda_d(K)\le3^d\sum_j\lambda_d(B_{i_j})<3^dk\sum_j\nu_s(B_{i_j})\le3^dk\,\nu_s(U)<3^dk\varepsilon.
\]
The union $\bigcup_xB_x$ is open. Inner regularity makes its measure the supremum of the measures of its compact subsets, so that measure is at most $3^dk\varepsilon$. Hence $\lambda_d^*(F_k)\le3^dk\varepsilon$ for every $\varepsilon$, and $F_k$ is null. The set $Z$ is null as well, so the limsup vanishes almost everywhere.
\end{proof}

\begin{think}
Can a continuous increasing function rise from $0$ to $1$ while its derivative vanishes almost everywhere?
\end{think}

Yes. The devil's staircase $c$ of Exercise~\ref{ex:devil} is constant on each open interval removed in the construction of the Cantor set, so $c'=0$ on a set of measure $1$, while $c(1)-c(0)=1$. The Stieltjes measure $\mu_c$ lives on the Cantor set. It is singular, and Theorem~\ref{thm:diffmeasures} sees only the density, which is $0$. The whole climb sits on a set of length $0$.

\begin{theorem}[Lebesgue, 1904]\label{thm:monotone}
Every monotone function $F:\R\to\R$ is differentiable at almost every point. If $F$ is non-decreasing and right-continuous, then $F'=\frac{d\mu_{F,a}}{d\lambda}$ a.e., where $\mu_{F,a}$ is the absolutely continuous part of
$\mu_F$ (Theorem~\ref{thm:stieltjes}), and $\int_a^bF'\le F(b)-F(a)$.
\end{theorem}
\begin{proof}
Let $F$ be non-decreasing and right-continuous, and write $\mu_F=f\lambda+\nu_s$. For $h>0$, the increment $F(x+h)-F(x)$ equals $\mu_F(]x,x+h])$, and $F(x)-F(x-h)$ equals $\mu_F(]x-h,x])$. Each of those intervals has length $h$ and sits inside $B(x,2h)$, which has length $4h$. At a point $x$ that is a Lebesgue point of $f$ and that satisfies $\nu_s(B(x,r))/(2r)\to0$, Theorems~\ref{thm:lebdiff} and~\ref{thm:diffmeasures} give
\[
\Bigl|\frac{F(x+h)-F(x)}h-f(x)\Bigr|\le\frac1h\int_x^{x+h}|f-f(x)|+\frac{\nu_s(B(x,2h))}h\to0,
\]
and the same limit from the left. The one-sided integral is at most the integral over $B(x,2h)$, whose average tends to $0$, and $\nu_s(B(x,2h))/h$ is four times the ratio $\nu_s(B(x,2h))/(4h)$, which tends to $0$. So $F'(x)=f(x)$ almost everywhere. Then
\[
\int_a^bF'=\int_{]a,b]}f\,d\lambda\le\mu_F(]a,b])=F(b)-F(a).
\]
The inequality can be strict. For the devil's staircase the integral of the derivative is $0$ and the increment is $1$.

Now let $F$ be an arbitrary non-decreasing function, and set $G(x)=F(x^+)$. Then $G$ is non-decreasing and right-continuous. The difference $H=G-F$ is nonnegative and vanishes off the discontinuities $x_j$ of $F$. There are only countably many discontinuities, because the jumps inside $[-n,n]$ are disjoint pieces of the rise $F(n+1)-F(-n-1)$. The measure $\rho=\sum_jH(x_j)\delta_{x_j}$ is locally finite and singular to $\lambda$. At every $x$ outside $\{x_j\}$ for which $\rho(B(x,r))/(2r)\to0$, one has $H(x)=0$ and
\[
0\le\frac{H(x+h)}{|h|}\le\frac{\rho(B(x,2|h|))}{|h|}\to0,
\]
since $H(x+h)$ is the jump at the single point $x+h$, if any. Theorem~\ref{thm:diffmeasures} puts almost every $x$ in that class. So $H'(x)=0$ and $F'=G'$ almost everywhere. A non-increasing function is the negative of a non-decreasing one.
\end{proof}

\section{Absolutely continuous functions}

\begin{definition}
$F:[a,b]\to\R$ is \emph{absolutely continuous} when, for every $\varepsilon>0$, some $\delta>0$ makes $\sum_i|F(b_i)-F(a_i)|<\varepsilon$ for every finite family of disjoint intervals $]a_i,b_i[$ of total length
$\sum_i(b_i-a_i)<\delta$.
\end{definition}

One interval makes this ordinary uniform continuity. Many intervals, taken together, rule out the devil's staircase: that function rises by $1$ on the Cantor set, and the Cantor set can be covered by open intervals of arbitrarily small total length.

\begin{theorem}[fundamental theorem of calculus]\label{thm:ftc}
For a function $F:[a,b]\to\R$, these are equivalent: (i) $F$ is absolutely continuous; (ii) some $f\in L^1([a,b])$ satisfies $F(x)=F(a)+\int_a^xf$ for all $x$. Then $F$ is differentiable a.e., $F'=f$ a.e., and
$F(b)-F(a)=\int_a^bF'$.
\end{theorem}
\begin{proof}
(ii)$\Rightarrow$(i). Lemma~\ref{lem:abscont} gives $\delta>0$ such that $\lambda(A)<\delta$ implies $\int_A|f|<\varepsilon$. For disjoint intervals of total length less than $\delta$, take $A$ to be their union. The sum of absolute increments is at most that integral. Differentiability, and $F'=f$ almost everywhere, are the one-dimensional case of Theorem~\ref{thm:lebdiff} recorded above.

(i)$\Rightarrow$(ii). \emph{Bounded variation.} Let $\delta$ be the number that absolute continuity attaches to $\varepsilon=1$. For a partition $P:a=x_0<\cdots<x_n=x$ write $V(F,P)=\sum_i|F(x_i)-F(x_{i-1})|$, and let $T(x)$ be the supremum of $V(F,P)$ over partitions of $[a,x]$. Cut $[a,b]$ into $m$ subintervals each of length less than $\delta$. On each of them every partition sum is $<1$, so the variation on each is at most $1$, and $T(b)\le m<\infty$. For $x<y$, concatenating a partition of $[a,x]$ with the single step from $x$ to $y$ gives $T(y)\ge T(x)+|F(y)-F(x)|$. So both $T$ and $T-F$ are non-decreasing. In fact $T(y)-T(x)$ equals the variation of $F$ on $[x,y]$. Concatenation gives the inequality $\ge$. For the other direction, add the point $x$ to any partition of $[a,y]$ and split the sum at $x$.

\emph{The variation is absolutely continuous.} Given $\varepsilon>0$, let $\delta$ be the number attached to this $\varepsilon$ for $F$. If $]a_i,b_i[$ are disjoint and have total length less than $\delta$, then any choice of partitions of the segments $[a_i,b_i]$ has total length less than $\delta$, so the sum of those partition sums is $<\varepsilon$. Taking suprema, $\sum_i\bigl(T(b_i)-T(a_i)\bigr)\le\varepsilon$. The definition asked for a strict inequality. Apply the same construction to $\varepsilon/2$: the sum of increments of $T$ is at most $\varepsilon/2$, hence strictly less than $\varepsilon$. So $T$ is absolutely continuous, and so is $T-F$. Absolute continuity with one interval is uniform continuity, so both functions are continuous.

\emph{A non-decreasing absolutely continuous function.} Let $G$ be non-decreasing and absolutely continuous on $[a,b]$, and extend it by the constant $G(a)$ to the left of $a$ and by the constant $G(b)$ to the right of $b$. The Stieltjes measure $\mu_G$ is absolutely continuous with respect to $\lambda$. Indeed, let $\lambda(N)=0$ and $\varepsilon>0$, and cover $N\cap\,]a,b[$ by open intervals of total length less than the $\delta$ attached to $\varepsilon$. On the line, a union of open intervals is a countable disjoint union of open intervals, of no larger total length, and the cover may be taken inside $]a,b[$. Continuity of $G$ gives $\mu_G(]a_i,b_i[)=G(b_i)-G(a_i)$. Every finite subcollection has sum of increments $<\varepsilon$, so $\mu_G(N)\le\varepsilon$. Since $\varepsilon$ is arbitrary, $\mu_G(N)=0$. The endpoints $a$ and $b$ are null for $\mu_G$ because $G$ is continuous. Lebesgue measure on $\R$ is $\sigma$-finite and $\mu_G$ is finite, so Theorem~\ref{thm:rn} supplies a density $g$ with $\mu_G=g\lambda$. Therefore $G(x)-G(a)=\mu_G(]a,x])=\int_a^xg$. Apply this to $T$ and to $T-F$. The difference of the two densities is the integrand required for $F$.
\end{proof}

The identity function $F(x)=x$ on $[0,1]$ is absolutely continuous, with integrand $1$, and $F(1)-F(0)=\int_0^11$. The devil's staircase is continuous, non-decreasing, and differentiable almost everywhere, with integrable derivative $0$, and it is not absolutely continuous. Absolute continuity is the condition the fundamental theorem uses, and the two directions of Theorem~\ref{thm:ftc} say that the condition is also sufficient.

\begin{beauty}
A density, a derivative, and a limit of averages over shrinking balls are three names for the same object, $\frac{d\nu}{d\lambda}$. The estimate that makes the limit exist almost everywhere is the batsman's most flattering average.
\end{beauty}

\section{Conditional expectation}

The number $\mathbb E[X]$ is the best constant guess of an integrable random variable $X$. Often something is known before the guess is made: the die came up even, the patient is in a given age band, the path up to time $s$ has been seen. The updated guess should be constant on each block of that information, and it should have the right average inside the block. The updated guess is what will be called the conditional expectation.

The case you can compute by hand is a finite partition. Let $A_1,\dots,A_n$ be disjoint events of positive probability whose union is $\Omega$, and let $\mathcal G$ be the $\sigma$-algebra they generate. A set lies in $\mathcal G$ when it is a union of some of the $A_i$. On the block $A_i$ the updated guess is the average of $X$ over that block,
\[
Y=\sum_{i=1}^n\frac{\mathbb E[X\mathbf 1_{A_i}]}{\mathbb P(A_i)}\,\mathbf 1_{A_i}.
\]
The random variable $Y$ is $\mathcal G$-measurable, because it is constant on each atom of $\mathcal G$. If $A$ is a union of some of the $A_i$, multiplication by $\mathbf 1_A$ keeps only those terms, so $\mathbb E[Y\mathbf 1_A]=\mathbb E[X\mathbf 1_A]$. Any other $\mathcal G$-measurable random variable with the same integrals on the atoms must take the same constant on each $A_i$.

\begin{example}[a fair die]\label{ex:die}
Let $X$ be the face shown by a fair die, and let $\mathcal G$ be the $\sigma$-algebra generated by the odd faces and the even faces. Then
\[
\frac{\mathbb E[X\mathbf 1_{\mathrm{odd}}]}{\mathbb P(\mathrm{odd})}=\frac{(1+3+5)/6}{1/2}=3,\qquad
\frac{\mathbb E[X\mathbf 1_{\mathrm{even}}]}{\mathbb P(\mathrm{even})}=\frac{(2+4+6)/6}{1/2}=4.
\]
So the updated guess $Y$ equals $3$ on the odd faces and $4$ on the even faces (Figure~\ref{fig:partition}). Averaging the two guesses recovers the ordinary expectation: $\frac12\cdot3+\frac12\cdot4=\frac72=\mathbb E[X]$. After Theorem~\ref{thm:condexp}, this random variable $Y$ is $\mathbb E[X\mid\mathcal G]$.
\end{example}

\begin{figure}[ht]
\centering
\fig{ch10-partition}
\caption{The conditional expectation of a fair die given the parity of the face. Each block carries the average of the faces inside it: $3$ on the three odd faces, drawn together, and $4$ on the three even faces. The guess is constant on the information that was given, and the average of the guess on each block equals the average of the face on that block.}
\label{fig:partition}
\end{figure}

Radon--Nikodym produces the same object for an arbitrary sub-$\sigma$-algebra. Form the measure that records the mass of $X$ on the sets one is allowed to see. Its density with respect to $\mathbb P$, computed on that smaller $\sigma$-algebra, is forced to be measurable with respect to it.

\begin{theorem}[conditional expectation]\label{thm:condexp}
Take $X\in L^1(\Omega,\mathcal F,\mathbb P)$, and let $\mathcal G\subseteq\mathcal F$ be a sub-$\sigma$-algebra. There is an integrable $\mathcal G$-measurable $Y$, written $\mathbb E[X\mid\mathcal G]$ and called the \emph{conditional expectation} of $X$ given $\mathcal G$, unique up to a null set, such that
\[
\mathbb E[Y\mathbf 1_A]=\mathbb E[X\mathbf 1_A]\qquad\text{for every }A\in\mathcal G.
\]
The map $X\mapsto\mathbb E[X\mid\mathcal G]$ is linear, $\mathbb E[\mathbb E[X\mid\mathcal G]]=\mathbb E[X]$, and $\mathbb E[ZX\mid\mathcal G]=Z\,\mathbb E[X\mid\mathcal G]$ whenever $Z$ is bounded and $\mathcal G$-measurable. If $X\ge0$ almost surely then $\mathbb E[X\mid\mathcal G]\ge0$ almost surely, and in general $|\mathbb E[X\mid\mathcal G]|\le\mathbb E[|X|\mid\mathcal G]$.
\end{theorem}
\begin{proof}
Suppose first that $X\ge0$. On the measurable space $(\Omega,\mathcal G)$, the set function $\nu(A)=\mathbb E[X\mathbf 1_A]$ is a finite measure. Countable additivity is monotone convergence applied to the partial sums of $X\mathbf 1_{A_k}$ (Theorem~\ref{thm:mct6}). If $\mathbb P(A)=0$, then $X\mathbf 1_A=0$ almost surely, so $\nu(A)=0$ by Proposition~\ref{prop:markov}(b). The restriction of $\mathbb P$ to $\mathcal G$ is a finite measure, hence $\sigma$-finite. Theorem~\ref{thm:rn}, applied on $(\Omega,\mathcal G,\mathbb P|_{\mathcal G})$, returns a $\mathcal G$-measurable $Y\ge0$ with $\mathbb E[Y\mathbf 1_A]=\nu(A)$ for every $A\in\mathcal G$. Taking $A=\Omega$ gives $\mathbb E[Y]=\mathbb E[X]<\infty$, so $Y$ is integrable and finite almost surely. For a general integrable $X$, both $X^+$ and $X^-$ are integrable and nonnegative. Set $\mathbb E[X\mid\mathcal G]=\mathbb E[X^+\mid\mathcal G]-\mathbb E[X^-\mid\mathcal G]$. The difference is defined almost surely because both terms are finite almost surely, and it has the integral identity of $X$.

If $Y$ and $Y'$ both satisfy the identity, the set $A=\{Y>Y'\}$ lies in $\mathcal G$ and $\mathbb E[(Y-Y')\mathbf 1_A]=0$ with $Y-Y'>0$ on $A$. Proposition~\ref{prop:markov}(b) forces $\mathbb P(A)=0$. Linearity follows from this uniqueness: the random variable $a\mathbb E[X\mid\mathcal G]+b\mathbb E[X'\mid\mathcal G]$ satisfies the integral identity that defines $\mathbb E[aX+bX'\mid\mathcal G]$. The case $A=\Omega$ is $\mathbb E[\mathbb E[X\mid\mathcal G]]=\mathbb E[X]$.

For the product rule, begin with $Z=\mathbf 1_B$ and $B\in\mathcal G$. For every $A\in\mathcal G$,
\[
\mathbb E[\mathbf 1_BY\mathbf 1_A]=\mathbb E[Y\mathbf 1_{A\cap B}]=\mathbb E[X\mathbf 1_{A\cap B}]=\mathbb E[\mathbf 1_BX\mathbf 1_A],
\]
so $\mathbf 1_BY=\mathbb E[\mathbf 1_BX\mid\mathcal G]$. Addition extends the identity to every simple $\mathcal G$-measurable $Z$. A bounded $\mathcal G$-measurable $Z$ is a uniform limit of simple $\mathcal G$-measurable functions $Z_n$, all bounded by the same constant $C$. Dominated convergence (Theorem~\ref{thm:dct6}), with the integrable bounds $C|X|$ and $C|Y|$, passes to the limit in $\mathbb E[Z_nY\mathbf 1_A]=\mathbb E[Z_nX\mathbf 1_A]$.

If $X\ge0$ at every point, the density produced by Theorem~\ref{thm:rn} is nonnegative. If $X\ge0$ only almost surely, pass to $\max(X,0)$. That function is nonnegative everywhere and equals $X$ almost surely, so the two have the same integrals on every set of $\mathcal G$. The nonnegative density built for $\max(X,0)$ therefore satisfies the identity for $X$. Uniqueness says it equals $\mathbb E[X\mid\mathcal G]$ almost surely. For a general integrable $X$, both $|X|+X$ and $|X|-X$ are nonnegative at every point, so $\mathbb E[|X|\mid\mathcal G]+\mathbb E[X\mid\mathcal G]\ge0$ and $\mathbb E[|X|\mid\mathcal G]-\mathbb E[X\mid\mathcal G]\ge0$ almost surely. Therefore $|\mathbb E[X\mid\mathcal G]|\le\mathbb E[|X|\mid\mathcal G]$.
\end{proof}

Conditioning in two stages does not depend on the order, as long as the second stage sees less than the first. Average $X$ on the fine partition, then average that result on the coarse partition. The result is the average of $X$ on the coarse partition.

\begin{theorem}[tower property]\label{thm:tower}
Take $X\in L^1(\Omega,\mathcal F,\mathbb P)$ and nested sub-$\sigma$-algebras $\mathcal H\subseteq\mathcal G\subseteq\mathcal F$. Then
\[
\mathbb E[\mathbb E[X\mid\mathcal G]\mid\mathcal H]=\mathbb E[X\mid\mathcal H]
\]
almost surely.
\end{theorem}
\begin{proof}
Set $Y=\mathbb E[X\mid\mathcal G]$. Theorem~\ref{thm:condexp} says that $\mathbb E[Y\mid\mathcal H]$ is $\mathcal H$-measurable and integrable. Every $A\in\mathcal H$ also lies in $\mathcal G$, so
\[
\mathbb E\bigl[\mathbb E[Y\mid\mathcal H]\,\mathbf 1_A\bigr]=\mathbb E[Y\mathbf 1_A]=\mathbb E[X\mathbf 1_A]=\mathbb E\bigl[\mathbb E[X\mid\mathcal H]\,\mathbf 1_A\bigr].
\]
The two $\mathcal H$-measurable random variables $\mathbb E[Y\mid\mathcal H]$ and $\mathbb E[X\mid\mathcal H]$ have the same integrals on every set in $\mathcal H$. Uniqueness in Theorem~\ref{thm:condexp} gives the identity.
\end{proof}

On the die of Example~\ref{ex:die}, take $\mathcal H=\{\varnothing,\Omega\}$. Then $\mathbb E[\mathbb E[X\mid\mathcal G]\mid\mathcal H]$ is the constant $\frac72$, which is also $\mathbb E[X\mid\mathcal H]$. Exercise~\ref{ex:condexp} runs the same arithmetic on a chain of two partitions that both split the die.

\section{What we proved, and what comes next}

A finite signed measure splits into a positive part and a negative part (Theorem~\ref{thm:hahn}). From that splitting, a $\sigma$-finite measure has a density with respect to another $\sigma$-finite measure exactly when it ignores the sets the other one ignores (Theorem~\ref{thm:rn}). Every $\sigma$-finite measure is a density part plus a singular part (Theorem~\ref{thm:lebdecomp}), and every law on the line has three parts (Corollary~\ref{cor:threeparts}). On $\R^d$ the density is a limit of averages over small balls. The maximal inequality controls those averages, and Lebesgue's differentiation theorem reads the limit (Theorems~\ref{thm:maximal}, \ref{thm:lebdiff} and~\ref{thm:diffmeasures}). Monotone functions are differentiable almost everywhere (Theorem~\ref{thm:monotone}). The fundamental theorem holds exactly for the absolutely continuous functions (Theorem~\ref{thm:ftc}). Conditional expectation exists, is unique up to a null set, and satisfies the tower property (Theorems~\ref{thm:condexp} and~\ref{thm:tower}). Each of these arguments is proved in this chapter, from the measure and the integral already built. Later volumes call on the conditional expectation and the tower property whenever they condition.

Every measure above lived on a single space. Probability also needs convergent sequences of laws, a law on a space of infinite sequences or of paths, and a measure of sets too thin for Lebesgue measure, such as a curve or the Cantor set. Those questions need a distance and a measure on the same space. Chapter~\ref{ch:metricmeasures} puts the two together.

\section{Exercises}

\begin{exercise}[$\star$]\label{ex:cantorfn}
Let $c$ be the devil's staircase of Exercise~\ref{ex:devil}. Show that $c'=0$ on $[0,1]\setminus C$, that $c$ is not absolutely continuous, and that its Lebesgue--Stieltjes measure $\mu_c$ is the Cantor law: singular, and without atoms. On $[0,\frac13]$, express the distribution function of the restriction of $\mu_c$ in terms of $c$.
\end{exercise}

\begin{exercise}[$\star$]\label{ex:condexp}
\emph{The tower, computed on a die.} Let $X$ be the face of a fair die. Let $\mathcal G$ be generated by $A_1=\{1,2\}$, $A_2=\{3,4\}$, $A_3=\{5,6\}$, and let $\mathcal H$ be generated by $B_1=A_1\cup A_2$ and $B_2=A_3$. Using the finite-partition formula, compute $Y=\mathbb E[X\mid\mathcal G]$ and $\mathbb E[X\mid\mathcal H]$, then compute $\mathbb E[Y\mid\mathcal H]$ the same way, and check that the two coarse guesses agree.
\end{exercise}

\begin{exercise}[$\star$]\label{ex:likelihood}
Let $\mathbb P=\mathcal N(0,1)$ and $\mathbb Q=\mathcal N(m,1)$ on $\R$. Compute $L=\frac{d\mathbb Q}{d\mathbb P}$, check $\mathbb E_{\mathbb P}[L]=1$, and compute the relative entropy $\int\ln L\,d\mathbb Q$. Show that $\mathcal N(0,1)$ and
$\mathcal N(0,2)$ are mutually absolutely continuous, and that $\mathcal N(0,1)$ and a Poisson law are mutually singular.
\end{exercise}

\begin{exercise}[$\star\star\star$]\label{ex:salem}
\emph{A strictly increasing singular function.} Let $0<p<1$, $p\ne\frac12$, let $B_1,B_2,\ldots$ be independent with $\mathbb P(B_k=1)=p=1-\mathbb P(B_k=0)$, and let $F_p$ be the distribution function of $X=\sum_kB_k2^{-k}$.
Show that $F_p$ is continuous and strictly increasing on $[0,1]$, and that $F_p'=0$ almost everywhere. Use Borel's theorem, Theorem~\ref{thm:borelnormal}, and prove the law of large numbers you need for these biased coins by the method of Lemma~\ref{lem:fourth}.
\end{exercise}

\begin{exercise}[$\star\star$]\label{ex:density}
\emph{Lebesgue's density theorem.} For a Borel set $A\subseteq\R^d$, show that $\frac{\lambda_d(A\cap B(x,r))}{\lambda_d(B(x,r))}\to\mathbf 1_A(x)$ for almost every $x$. Deduce that no Borel set $A\subseteq[0,1]$ satisfies
$\lambda(A\cap I)=\frac12\lambda(I)$ for every interval $I\subseteq[0,1]$. Compare the conclusion with Exercise~\ref{ex:everywhere}.
\end{exercise}

\begin{exercise}[$\star\star$]\label{ex:variation}
(a) Show that $x\sin\frac1x$, with value $0$ at $0$, is continuous on $[0,1]$ and is not of bounded variation. (b) Show that $F(x)=x^2\sin\frac1{x^2}$, with $F(0)=0$, is differentiable at every point of $[-1,1]$, and that
$F'$ is not integrable on $[0,1]$. Is $F$ absolutely continuous on $[0,1]$?
\end{exercise}

\begin{exercise}[$\star$]\label{ex:maximalcompute}
\emph{A maximal function, computed by hand.} Let $f=\mathbf 1_{[-1,1]}$ on $\R$. (a) Show that $Mf(x)=1$ for $|x|<1$ and $Mf(x)=\frac1{|x|+1}$ for $|x|\ge1$, as drawn in Figure~\ref{fig:maximal}.
(b) Compute $\lambda(\{Mf>t\})$ for $0<t<1$, and compare the value with the bound in Theorem~\ref{thm:maximal}. (c) Show that $Mf$ is not integrable.
\end{exercise}

\begin{exercise}[$\star$]\label{ex:v7c11a}
Let $\nu$ be the signed measure on the Borel sets of $\R$ given by $\nu(A)=\lambda(A\cap[0,2])-2\lambda(A\cap[-1,0])$. Find a Hahn partition, write $\nu^+$ and $\nu^-$, and compute $|\nu|(\R)$. If you move the half-line $(2,\infty)$ from one side of the partition to the other, what happens to $\nu$ on subsets of that half-line?
\end{exercise}

\begin{exercise}[$\star\star$]\label{ex:v7c11b}
Let $\mu$ be counting measure on the Borel sets of $\R$: $\mu(A)$ is the number of points of $A$ when $A$ is finite, and $\mu(A)=\infty$ otherwise. Chapter~\ref{ch:products} already uses counting measure on $[0,1]$. Show that Lebesgue measure satisfies $\lambda\ll\mu$, that $\mu$ is not $\sigma$-finite, and that no measurable $f\ge0$ satisfies $\lambda(A)=\int_Af\,d\mu$ for every Borel $A$.
\end{exercise}

\begin{exercise}[$\star\star$]\label{ex:v7c11c}
Let $\mu$, $\nu$ and $\lambda$ be finite measures with $\lambda\ll\nu\ll\mu$. Let $f$ be a finite nonnegative version of $\frac{d\nu}{d\mu}$, and let $g$ be a finite nonnegative version of $\frac{d\lambda}{d\nu}$. Show that $gf=\frac{d\lambda}{d\mu}$ holds $\mu$-almost everywhere.
\end{exercise}

\begin{exercise}[$\star$]\label{ex:v7c11d}
Let $X$ be integrable and $\mathcal G$-measurable. Show that $\mathbb E[X\mid\mathcal G]=X$ almost surely. If $\mathcal H=\{\varnothing,\Omega\}$, show that $\mathbb E[X\mid\mathcal H]$ equals the constant $\mathbb E[X]$. On the die of Example~\ref{ex:die}, check that the face $X$ is not $\mathcal G$-measurable, that $Y=\mathbb E[X\mid\mathcal G]$ satisfies $\mathbb E[Y\mid\mathcal G]=Y$, and that $\mathbb E[Y\mid\mathcal H]=\frac72$.
\end{exercise}

\section{Hints}

\begin{hint}{ex:cantorfn}
A continuous function that is constant on an open interval has derivative $0$ there. For absolute continuity, cover the stage $C_n$ by $2^n$ intervals whose lengths add to $(\frac23)^n$. On $[0,\frac13]$ the staircase satisfies $c(x)=\frac12c(3x)$.
\end{hint}
\begin{hint}{ex:condexp}
The average of the face on $A_1$ is $\frac32$. On the coarse block $B_1$, average the two values that $Y$ takes there. Each of $A_1$ and $A_2$ contains two of the four faces in $B_1$.
\end{hint}
\begin{hint}{ex:likelihood}
The likelihood ratio is the quotient of the two Gaussian densities. Under $\mathbb Q$ one has $\ln L(x)=mx-\frac{m^2}2$, and the mean of $X$ is $m$.
\end{hint}
\begin{hint}{ex:salem}
Every dyadic interval of rank $n$ has $\mu_p$-probability $p^k(1-p)^{n-k}>0$. The set of $x$ whose binary digits have frequency of ones equal to $p$ has $\mu_p$-measure $1$ and Lebesgue measure $0$. Theorem~\ref{thm:monotone} then forces the derivative to vanish.
\end{hint}
\begin{hint}{ex:density}
Apply Theorem~\ref{thm:lebdiff} to the integrable function $\mathbf 1_{A\cap B(0,R)}$, and let $R$ run through the positive integers. If every interval met $A$ in half its length, the density at almost every point of $[0,1]$ would equal $\frac12$.
\end{hint}
\begin{hint}{ex:variation}
(a) Test the variation at the points $x_k=\frac1{(k+\frac12)\pi}$. (b) For $x\ne0$, $F'(x)=2x\sin\frac1{x^2}-\frac2x\cos\frac1{x^2}$. Integrate $|F'|$ on the intervals where $\cos\frac1{x^2}\ge\frac12$.
\end{hint}
\begin{hint}{ex:maximalcompute}
(a) For $x>1$, the interval centred at $x$ meets $[-1,1]$ only when its radius exceeds $x-1$, and the average is largest for the smallest radius that contains the whole of $[-1,1]$.
\end{hint}
\begin{hint}{ex:v7c11a}
On $[0,2]$ only the first term of $\nu$ is alive, and it is Lebesgue measure. On $[-1,0]$ only the second term is alive, and the factor $2$ does not change its sign. Off $[-1,2]$ both terms are zero.
\end{hint}
\begin{hint}{ex:v7c11b}
A set of counting measure zero is empty, so absolute continuity is immediate. A singleton has Lebesgue measure zero, which forces $f(x)=0$ at every $x$. The unit interval then has the wrong measure. A finite counting-measure set is a finite set, and $\R$ is uncountable.
\end{hint}
\begin{hint}{ex:v7c11c}
For an indicator, integrating $g$ against $\nu$ returns $\lambda(A)$. Replacing $\nu$ by $f\mu$ turns that integral into an integral of $gf$ against $\mu$. Uniqueness in Theorem~\ref{thm:rn} names the product.
\end{hint}
\begin{hint}{ex:v7c11d}
If $X$ is already $\mathcal G$-measurable, it satisfies the integral identity that defines the conditional expectation. On the trivial $\sigma$-algebra the measurable random variables are the constants, and the identity on $\Omega$ fixes the constant. On the die, $Y$ is constant on the odd faces and constant on the even faces.
\end{hint}

\section{Solutions}

\begin{solution}{ex:cantorfn}
\idea The staircase is flat on the gaps of the Cantor set, and those gaps have total length $1$. It climbs only on the Cantor set, which open intervals of tiny total length can cover. That cover is what absolute continuity forbids.
\writeup Exercise~\ref{ex:devil} makes $c$ constant on each removed open interval, so $c'=0$ on $[0,1]\setminus C$, a set of measure $1$. The stage $C_n$ is $2^n$ closed intervals of length $3^{-n}$. Self-similarity gives $c(x)=\frac12c(3x)$ on $[0,\frac13]$ and $c(x)=\frac12+\frac12c(3x-2)$ on $[\frac23,1]$, so $c$ rises by $2^{-n}$ on each interval of stage $n$. The lengths add to $(\frac23)^n$, which tends to $0$, while the sum of the rises is $1$. Enlarge the intervals slightly to open intervals of total length less than any prescribed $\delta$. Then $c$ fails absolute continuity. Extend $c$ by $0$ on $]-\infty,0[$ and by $1$ on $]1,\infty[$. The measure $\mu_c$ gives mass $0$ to each gap, because $c$ is constant there, so $\mu_c([0,1]\setminus C)=0$. Also $\lambda(C)=0$. Continuity of $c$ forces $\mu_c(\{x\})=0$. So $\mu_c$ is singular and atomless. It is the law of $\sum_k2B_k3^{-k}$ with fair bits $B_k$, because both measures have distribution function $c$, again by self-similarity. Restricted to $[0,\frac13]$, the distribution function is $c(x)=\frac12c(3x)$: half the image of $\mu_c$ under $x\mapsto\frac x3$.
\end{solution}

\begin{solution}{ex:condexp}
\idea Each block contributes its own average. The tower says that averaging the fine averages over a coarse block is the same arithmetic as averaging the face over that block directly.
\writeup The finite-partition formula gives $Y=\frac32$ on $A_1$, $Y=\frac72$ on $A_2$, and $Y=\frac{11}2$ on $A_3$. On $B_1=A_1\cup A_2$,
\[
\mathbb E[X\mid\mathcal H]=\frac{1+2+3+4}{4}=\frac52,\qquad
\mathbb E[Y\mid\mathcal H]=\frac{\frac32+\frac32+\frac72+\frac72}{4}=\frac52,
\]
because $A_1$ and $A_2$ each contain two of the four faces in $B_1$. On $B_2=A_3$ both conditional expectations equal $\frac{11}2$. Therefore $\mathbb E[Y\mid\mathcal H]=\mathbb E[X\mid\mathcal H]$. This is Theorem~\ref{thm:tower} for the chain $\mathcal H\subseteq\mathcal G$.
\end{solution}

\begin{solution}{ex:likelihood}
\idea Both Gaussian laws have densities that are positive at every real, so each is absolutely continuous with respect to the other, and the likelihood ratio is the quotient of those densities. A Poisson law lives on the integers.
\writeup Write $\varphi(x)=\frac{e^{-x^2/2}}{\sqrt{2\pi}}$. Then $\mathbb Q(A)=\int_A\varphi(x-m)\,dx=\int_A\frac{\varphi(x-m)}{\varphi(x)}\,d\mathbb P(x)$, so $L(x)=e^{mx-m^2/2}$. Hence $\mathbb E_{\mathbb P}[L]=\mathbb Q(\R)=1$. The same identity is the Gaussian moment $\mathbb E[e^{mZ}]=e^{m^2/2}$. Next,
\[
\int\ln L\,d\mathbb Q=\mathbb E_{\mathbb Q}[mX]-\frac{m^2}2=m\cdot m-\frac{m^2}2=\frac{m^2}2.
\]
The laws $\mathcal N(0,1)$ and $\mathcal N(0,2)$ both have strictly positive densities, so the same quotient runs in both directions and the two measures are mutually absolutely continuous. A Poisson law is carried by $\N$, and $\mathcal N(0,1)$ gives $\N$ measure $0$, so the two laws are mutually singular.
\end{solution}

\begin{solution}{ex:salem}
\idea Under $\mu_p$ the binary digits of $X$ are biased coins, so the law of large numbers puts the proportion of ones at $p$. Under Lebesgue measure that proportion is $\frac12$. The two measures live on disjoint sets. Every dyadic interval still has positive $\mu_p$-mass, so the distribution function rises on every interval.
\writeup Let $\mu_p$ be the law of $X$. A single point has at most two binary expansions, and each expansion has probability at most $\lim_n\max(p,1-p)^n=0$. So $\mu_p$ has no atoms, and $F_p$ is continuous. The dyadic interval $[\frac{j}{2^n},\frac{j+1}{2^n}]$ contains $X$ as soon as the first $n$ bits are those of $j$, an event of probability $p^k(1-p)^{n-k}>0$. Every nontrivial subinterval of $[0,1]$ contains such a dyadic interval, so $F_p(t)>F_p(s)$ whenever $s<t$.

For the law of large numbers, set $Y_k=B_k-p$. Then $|Y_k|\le1$ and $\mathbb E[Y_k]=0$. Expand $\mathbb E[(Y_1+\cdots+Y_n)^4]$ by the same bookkeeping as Lemma~\ref{lem:fourth}. The coins are independent, so the expectation of a product of bounded functions of distinct coins is the product of the expectations (Chapter~\ref{ch:products}, applied to two factors and then grouped). Every term in which some index appears exactly once contains a factor $\mathbb E[Y_i]=0$ and vanishes. What remains is $n$ terms $\mathbb E[Y_k^4]\le1$, together with $3n(n-1)$ terms $\mathbb E[Y_j^2]\mathbb E[Y_k^2]\le1$. The whole expectation is at most $3n^2$. Markov's inequality gives
\[
\mathbb P\Bigl(\Bigl|\sum_{k\le n}Y_k\Bigr|\ge n\varepsilon\Bigr)\le\frac{3}{\varepsilon^4n^2}.
\]
The right-hand side is summable in $n$. Borel--Cantelli, Proposition~\ref{prop:measureprops}(e), yields $\frac1n\sum_{k\le n}B_k\to p$ almost surely.

Let $S_q$ be the set of $x\in[0,1[$ whose binary digits $D_j(x)$ satisfy $\frac1n\sum_{j\le n}D_j(x)\to q$. The law just proved gives $\mu_p(S_p)=1$, except on the null event that $X$ is dyadic, where the expansion is not unique. Borel's theorem, Theorem~\ref{thm:borelnormal}, gives $\lambda(S_{1/2})=1$. The sets $S_p$ and $S_{1/2}$ are disjoint because $p\ne\frac12$, so $\lambda(S_p)=0$. Thus $\mu_p\perp\lambda$. The absolutely continuous part is $0$, and Theorem~\ref{thm:monotone} forces $F_p'=0$ almost everywhere. The function is continuous and strictly increasing, and the climb sits on a dense set of measure zero.
\end{solution}

\begin{solution}{ex:density}
\idea The proportion of $A$ inside a small ball is the average of $\mathbf 1_A$ over that ball. Lebesgue's theorem says that those averages return the value of the function at almost every point, and $\mathbf 1_A$ takes only the values $0$ and $1$.
\writeup Fix $R$ and set $f=\mathbf 1_{A\cap B(0,R)}$, which is integrable. Theorem~\ref{thm:lebdiff} says that for almost every $x$ with $|x|<R$, and for all sufficiently small $r$,
\[
\frac{\lambda_d(A\cap B(x,r))}{\lambda_d(B(x,r))}=\frac1{\lambda_d(B(x,r))}\int_{B(x,r)}f\to f(x)=\mathbf 1_A(x).
\]
Let $R$ run through the positive integers. If $\lambda(A\cap I)=\frac12\lambda(I)$ for every subinterval of $[0,1]$, the density of $A$ at every point of $]0,1[$ equals $\frac12$, which is neither $0$ nor $1$. The set where the conclusion fails is then all of $]0,1[$, which is not null. Exercise~\ref{ex:everywhere} built a set that meets every interval in positive measure, and whose complement does the same. The density theorem allows that. It forbids the proportions from being exactly one half in every interval.
\end{solution}

\begin{solution}{ex:variation}
\idea Bounded variation cannot afford infinitely many wiggles of height about $\frac1k$, because the harmonic series diverges. A derivative can be large enough near a single point to leave $L^1$, and then Theorem~\ref{thm:ftc} refuses absolute continuity.
\writeup (a) Continuity at $0$ is $|x\sin\frac1x|\le|x|$. At $x_k=\frac1{(k+\frac12)\pi}$ one has $x_k\sin\frac1{x_k}=\frac{(-1)^k}{(k+\frac12)\pi}$, so consecutive sample values differ by more than $\frac1{(k+1)\pi}$. The variation of the function on $[0,1]$ is at least $\sum_k\frac1{(k+1)\pi}=\infty$.

(b) The formula in the hint holds for $x\ne0$. At the origin, $\bigl|\frac{F(h)}h\bigr|\le|h|\to0$, so $F'(0)=0$. On the interval
\[
I_k=\bigl[(2k\pi+\tfrac\pi3)^{-1/2},\,(2k\pi-\tfrac\pi3)^{-1/2}\bigr]
\]
one has $\cos\frac1{x^2}\ge\frac12$. For $x\le\frac12$, $|F'(x)|\ge\frac1x-2x\ge\frac1{2x}$. Therefore
\[
\int_{I_k}|F'|\ge\frac12\ln\frac{(2k\pi+\frac\pi3)^{1/2}}{(2k\pi-\frac\pi3)^{1/2}}=\frac14\ln\frac{6k+1}{6k-1}\ge\frac c k
\]
for some $c>0$ and all large $k$. The series $\sum\frac ck$ diverges, so $F'$ is not integrable on $[0,1]$. Theorem~\ref{thm:ftc} says that an absolutely continuous function on $[0,1]$ would have an integrable derivative. So $F$ is not absolutely continuous on $[0,1]$. On each $[\varepsilon,1]$ the derivative is bounded, and $F$ is absolutely continuous there. The failure is concentrated at the origin.
\end{solution}

\begin{solution}{ex:maximalcompute}
\idea Far from $[-1,1]$, the most flattering interval centred at $x$ is the shortest one that still contains the whole of $[-1,1]$. A shorter interval misses some of the mass. A longer one only adds empty length.
\writeup (a) If $|x|<1$, a small interval about $x$ lies inside $[-1,1]$, and the average equals $1$, which is the largest value $|f|$ ever takes. By symmetry let $x\ge1$. The interval $]x-r,x+r[$ has length $2r$. Its intersection with $[-1,1]$ has length $0$ when $r\le x-1$, length $r-(x-1)$ when $x-1\le r\le x+1$, and length $2$ when $r\ge x+1$. On the middle range the average is $\frac{r-x+1}{2r}=\frac12-\frac{x-1}{2r}$, which increases with $r$. On the last range the average is $\frac1r$, which decreases with $r$. The supremum is attained at $r=x+1$, and $Mf(x)=\frac1{x+1}$. At the endpoint $x=1$ this value is $\frac12$: half of every small interval about the endpoint falls outside $[-1,1]$.

(b) For $0<t<1$, the inequality $Mf(x)>t$ holds for every $|x|<1$, and it holds for $|x|\ge1$ precisely when $\frac1{|x|+1}>t$, that is $|x|<\frac1t-1$. If $t\le\frac12$, then $\frac1t-1\ge1$, so the set is $|x|<\frac1t-1$ and its length is $\frac2t-2$. If $\frac12<t<1$, then $\frac1{|x|+1}\le\frac12<t$ for every $|x|\ge1$, while $Mf=1>t$ on $|x|<1$, so the length is $2$. The bound in Theorem~\ref{thm:maximal} is $\frac3t\|f\|_1=\frac6t$. It has the same order $\frac1t$, and the constant in front of $\frac1t$ is three times the constant in $\frac2t$.

(c) $\int_1^R\frac{dx}{x+1}=\ln\frac{R+1}2$, which tends to infinity with $R$. So $Mf\notin L^1$, even though $f$ is bounded and has bounded support. An inequality $\|Mf\|_1\le C\|f\|_1$ is impossible for a finite $C$. The bound on $\lambda(\{Mf>t\})$ is the statement that remains true.
\end{solution}

\begin{solution}{ex:v7c11a}
\idea The first term of $\nu$ sees only $[0,2]$ and is nonnegative there. The second term sees only $[-1,0]$ and is nonpositive there. Off $[-1,2]$ the signed measure is zero, so that region may join either side of a Hahn partition.
\writeup Take $P=[0,2]$ and $N=\R\setminus[0,2]$. If $A\subseteq P$, then $\nu(A)=\lambda(A)\ge0$, so $P$ is positive. If $A\subseteq N$, then $A$ misses $[0,2]$ and $\nu(A)=-2\lambda(A\cap[-1,0])\le0$, so $N$ is negative. Jordan's formulas give $\nu^+(A)=\lambda(A\cap[0,2])$ and $\nu^-(A)=2\lambda(A\cap[-1,0])$. Hence
\[
|\nu|(\R)=\nu([0,2])-\nu([-1,0])=2-(-2)=4.
\]
The half-line $(2,\infty)$ satisfies $\nu(B)=0$ for every measurable $B\subseteq(2,\infty)$. Moving it into $P$ keeps $P$ positive, because a subset of $[0,2]\cup(2,\infty)$ still meets $\nu$ only through Lebesgue measure on $[0,2]$. This is the uniqueness clause of Theorem~\ref{thm:hahn}: the symmetric difference of the two positive sets is a set on which $\nu$ vanishes.
\end{solution}

\begin{solution}{ex:v7c11b}
\idea Absolute continuity only has to be checked on the empty set, because that is the only set of counting measure zero. The Radon--Nikodym derivative would have to reproduce the length of every singleton, which is zero, and also the length of $[0,1]$, which is one. Those two demands cannot be met by one function. The real line is too large to be a countable union of finite sets.
\writeup The empty set is the only set with $\mu(A)=0$. Lebesgue measure of the empty set is $0$, so $\lambda\ll\mu$. A set of finite counting measure is a finite set. A countable union of finite sets is countable. Theorem~\ref{thm:continuum} says that $\R$ is uncountable, so $\mu$ is not $\sigma$-finite. Suppose a measurable $f\ge0$ satisfied $\lambda(A)=\int_Af\,d\mu$ for every Borel $A$. For each $x$, the singleton $\{x\}$ is Borel and $\lambda(\{x\})=0$. If $f(x)=\infty$, the integral of $f$ over $\{x\}$ would be infinite. If $f(x)$ is finite, that integral equals $f(x)$. Either way $f(x)=0$. The zero function has integral $0$, by the convention $0\cdot\infty=0$ in the definition of the integral of a simple function. But $\lambda([0,1])=1$. So no such $f$ exists, and $\sigma$-finiteness cannot be omitted from Theorem~\ref{thm:rn}.
\end{solution}

\begin{solution}{ex:v7c11c}
\idea The derivative $g$ turns integrals against $\nu$ into integrals against $\lambda$. The derivative $f$ turns integrals against $\mu$ into integrals against $\nu$. Composing the two substitutions multiplies the densities, and uniqueness names the product.
\writeup Integrability of $g$ with respect to $\nu$ forces $g<\infty$ for $\nu$-almost every point, by Proposition~\ref{prop:markov}(c). Let $Z=\{g=\infty\}$. Then $\nu(Z)=0$, so $\int_Zf\,d\mu=0$, and Proposition~\ref{prop:markov}(b) gives $f=0$ for $\mu$-almost every point of $Z$. Redefine $g$ to be $0$ on $Z$. The new $g$ is finite, nonnegative, and still satisfies $\int_Ag\,d\nu=\lambda(A)$ for every measurable $A$, because $Z$ is $\nu$-null. For $g=\mathbf 1_B$,
\[
\int_A\mathbf 1_Bf\,d\mu=\int_{A\cap B}f\,d\mu=\nu(A\cap B)=\int_A\mathbf 1_B\,d\nu.
\]
Nonnegative simple functions follow by addition, Corollary~\ref{cor:linearity6}. For a general measurable $g\ge0$, choose simple $g_n$ increasing to $g$. Then $g_nf$ increases to $gf$, and monotone convergence (Theorem~\ref{thm:mct6}) passes to the limit on both sides. Therefore $\int_A(gf)\,d\mu=\int_Ag\,d\nu=\lambda(A)$ for every measurable $A$. Theorem~\ref{thm:rn} says that the density is unique $\mu$-almost everywhere, so $gf=\frac{d\lambda}{d\mu}$ holds $\mu$-almost everywhere.
\end{solution}

\begin{solution}{ex:v7c11d}
\idea The defining identity asks for a $\mathcal G$-measurable random variable with the same integrals as $X$ on the sets in $\mathcal G$. If $X$ itself is $\mathcal G$-measurable, it is already that random variable. On the trivial $\sigma$-algebra the only measurable functions are constants, and the integral over $\Omega$ picks out $\mathbb E[X]$.
\writeup Suppose $X\ge0$ is integrable and $\mathcal G$-measurable. Then $\mathbb E[X\mathbf 1_A]=\mathbb E[X\mathbf 1_A]$ for every $A\in\mathcal G$, so $X$ satisfies the identity in Theorem~\ref{thm:condexp}. Uniqueness gives $\mathbb E[X\mid\mathcal G]=X$ almost surely. For a general integrable $\mathcal G$-measurable $X$, the maps $t\mapsto\max(t,0)$ and $t\mapsto\max(-t,0)$ are continuous, so $X^+$ and $X^-$ are $\mathcal G$-measurable (Chapter~\ref{ch:integral}). Both are integrable. The nonnegative case yields $\mathbb E[X^+\mid\mathcal G]=X^+$ and $\mathbb E[X^-\mid\mathcal G]=X^-$ almost surely, and their difference is $X$.

A random variable is $\mathcal H$-measurable, for $\mathcal H=\{\varnothing,\Omega\}$, precisely when it is constant almost surely. The constant $c=\mathbb E[X]$ satisfies $\mathbb E[c\mathbf 1_\Omega]=\mathbb E[X]$ and $\mathbb E[c\mathbf 1_\varnothing]=0$. Uniqueness says that this constant is $\mathbb E[X\mid\mathcal H]$.

On the die, $\mathcal G$ is generated by the odd faces and the even faces. The face $X$ takes three different values on the odd faces, so it is not constant on the atoms of $\mathcal G$ and is not $\mathcal G$-measurable. The conditional expectation $Y$ equals $3$ on the odds and $4$ on the evens, so it is $\mathcal G$-measurable, and the first part gives $\mathbb E[Y\mid\mathcal G]=Y$. Example~\ref{ex:die} computed $\mathbb E[Y]=\frac72$. The second part says that $\mathbb E[Y\mid\mathcal H]$ is that constant.
\end{solution}

\section*{Sources and further reading}
\addcontentsline{toc}{section}{Sources and further reading}

\begin{itemize}
\item The differentiation of monotone functions, and the fundamental theorem for the Lebesgue integral, are in H.~Lebesgue, \emph{Le\c cons sur l'int\'egration et la recherche des fonctions primitives}, Gauthier-Villars, 1904.
\item The Euclidean case is J.~Radon, ``Theorie und Anwendungen der absolut additiven Mengenfunktionen'', \emph{Sitzungsber. Akad. Wiss. Wien} 122 (1913), 1295--1438. The general measure space is O.~Nikodym, ``Sur une g\'en\'eralisation des int\'egrales de M.~J.~Radon'', \emph{Fund. Math.} 15 (1930), 131--179.
\item G.~H.~Hardy and J.~E.~Littlewood, ``A maximal theorem with function-theoretic applications'', \emph{Acta Math.} 54 (1930), 81--116. The article begins on page 81. The cricket comparison opens Section~II, on page 83.
\item G.~B.~Folland, \emph{Real Analysis}, 2nd ed., Wiley, 1999, Chapter~3, is the source of the route through the maximal function and the differentiation of measures.
\item A second account of the same theorems is E.~M.~Stein and R.~Shakarchi, \emph{Real Analysis}, Princeton, 2005, Chapter~3.
\item The strictly increasing singular function of Exercise~\ref{ex:salem} is the subject of R.~Salem, ``On some singular monotonic functions which are strictly increasing'', \emph{Trans. Amer. Math. Soc.} 53 (1943), 427--439.
\item The cricket paragraph is quoted from the Hardy--Littlewood paper above. The derivative that exists everywhere and fails to be Riemann integrable is V.~Volterra, ``Sui principii del calcolo integrale'', \emph{Giornale di Matematiche} 19 (1881), 333--372.
\end{itemize}
